\documentstyle[%


righttag%


,11pt%


,amssymb%


,russian%


,izvru%


]{amsart}





\newcommand{\tl}[3]{\medskip\noindent{\bf#1}\ #2 \dotfill #3 \par} 


\pagestyle{empty}


\begin{document}


%\tableofcontents


\par


\vskip 2cm


\par


\bigskip


\bigskip


\bigskip


\par


\centerline{\Large СОДЕРЖАНИЕ}


\bigskip


\bigskip


%%%%%%%%%% Use \newline %%%%%%%%%%%





\tl{Абдурагимов Э.И.}


{Единственность положительного решения 


одной нелинейной\newline двухточечной краевой задачи 


и численный метод его нахождения}{3}


\tl{Бадриев И.Б., Ляшко А.Д., Панкратова О.В.}


{Исследование сходимости\newline итерационных методов 


решения нелинейных задач теории фильтрации}{8}


\tl{Булатов М.В.} 


{Редукция вырожденных систем интегральных  


уравнений типа Вольтерра\newline к системам 2-го рода}{14}


\tl{Вербицкий В.В.}


{Смешанный метод конечных элементов 


в задаче на собственные\newline значения 


нелинейной устойчивости пологих оболочек}{22}


\tl{Калинин Д.А.}


{$H$-проективно-эквивалентные 


римановы связности}{32}


\tl{Каюмов И.Р.}


{Об одной экстремальной задаче теории крыла}{41}


\tl{Короткий А.И.} 


{Восстановление управлений и параметров 


динамических систем при\newline неполной информации}{47}


\tl{Кузнецов Е.Б., Шалашилин В.И.}


{Решение сингулярных уравнений, 


преобразованных\newline к наилучшему аргументу}{56}


\tl{Мельников Ю.Б.}


{Представление полугруппы эндоморфизмов 


конечной группы на\newline пространстве ее классовых функций}{64}


\tl{Павленко В.Н., Ульянова О.В.} 


{Метод верхних и нижних решений для уравнений\newline


эллиптического типа с разрывными нелинейностями}{69}


\tl{Перов А.И.} 


{Об одной теореме Островского}{77}


\tl{Солодкий С.Г.}


{О модификации проекционной схемы 


решения некорректных задач}{83}








\bigskip


{\centerline {Краткие сообщения}}


\bigskip





\tl{Абдуллаев Р.З.} 


{Критерий изоморфности  некоммутативных алгебр Аренса}{90}


\tl{Бутаков В.В., Мартынов Л.М.}


{Полугруппы, все нециклические подполугруппы 


которых\newline являются идеалами}{94}


\bigskip





\noindent{}Рефераты статей, депонированных в ВИНИТИ


через журнал ``Известия вузов.\newline Математика''\dotfill{97}











\vfill


\noindent\raisebox{2pt}{\copyright} Казанский госудаpственный унивеpситет


\end{document}


